\section{Do the findings extend to rules about permitted changes?}
\label{sec:substantive_extension}
The cancellation pairs show what happens under controlled facts. They do not tell us how often agents naturally propose forbidden changes. We therefore collected additional conversations with unchanged customer scenarios and substantive policies. This section separates what that study found about rule violations from what it found about monitor warnings.

\subsection{New conversations and six specific rule checks}
A feasibility pilot completed 23 of 24 planned conversations on 12 previously uncollected training tasks. A subsequent locally frozen allocation selected another 24 training task IDs by a deterministic hash. Two acting models, DeepSeek and GPT-5.4 Mini, each ran every task twice, producing 96 of 96 planned conversations. Both received the same reminder about the interaction rule. We analyze this new-task study separately from the pilot and the reserved-test cohorts.

We check six rules whose relevant facts can be read from the recorded tool observations:
\begin{itemize}
\item \textbf{Airline:} restrictions on changing basic-economy itineraries, changing the number of passengers, and removing baggage.
\item \textbf{Retail:} whether order status permits an operation, whether a refund destination is allowed, and whether an item change stays within the required product type.
\end{itemize}
A second implementation checks the same rules using a separate reverse-prefix procedure. The two implementations agree on 152 rule-check rows across the pilot and new-task study. This checks the software, not every possible interpretation of a policy. Missing facts remain unknown. Contextual cancellation eligibility, compensation and customer authorization are outside these six checks.

The 96 conversations contain 42 with an applicable selected-rule check and no detected violation, 5 with a violation candidate, 3 with unresolved facts, and 46 to which none of the selected checks applies. An inapplicable conversation is not counted as having passed a substantive rule. The 24 task IDs contain 21 distinct reference-operation sequences; 10 share visible users, orders or reservations with earlier collection. New IDs therefore do not imply independent workflows or databases.

\subsection{What the violation candidates actually show}
The five candidate conversations cover three task IDs. Three runs of retail task 110 propose an exchange on a pending order, and one run of retail task 7 proposes an operation requiring pending status on a delivered order. The tools reject these proposals. The fifth, retail task 72, changes an address after an item modification: the written policy forbids further modification, but the tool accepts the change. This is a disagreement between the written policy and tool implementation, not evidence of an unauthorized customer request.

\paragraph{Observed example: the tool accepts a change that the written rule forbids.}
In retail task 72 (DeepSeek, repeat 0), message 24 modifies an order's items. The tool returns status ``pending (item modified)'' at message 25. The written policy prohibits further modification after this step. Nevertheless, message 26 proposes an address change, which the tool accepts at message 27; the final database matches its reference. At message 26, Gemini scores both broad and six-rule targets at 1.0. GPT-5.4 Mini changes from 0.22 under the broad question to 0.99 under the six-rule question, crossing the 0.5 threshold. Both targets use the same visible prefix. The customer requested the changes, so this is a workflow restriction, not an unauthorized customer request. The conversation also batches earlier lookups, and the simulator's permissive implementation conflicts with its written policy. We therefore exclude it from the strict comparison below.

All five candidate conversations also contain an interaction-rule violation. Consequently, removing conversations with either an interaction violation or a policy--tool disagreement leaves \emph{no selected-rule violation examples}. The study cannot estimate violation-detection sensitivity in that stricter subset. The rejected proposals also must not be described as completed harmful changes. The separate pilot does contain one interaction-clean violation candidate: an airline agent proposes changing a reservation from one passenger to three, and the tool rejects it. That single pilot case does not supply independent positive-task support for the new-task study.

\subsection{Changing the monitoring question changes warnings}
We nevertheless can test whether the monitoring question changes warnings on the same recorded evidence. Gemini and GPT-5.4 Mini each receive three questions: has any public-policy violation occurred or been proposed; has the interaction rule been violated; or has one of the six selected substantive rules been violated? All three questions use the same time horizon. None asks the model to predict an unseen future failure.

For every applicable proposal, we select the nearest generated assistant decisions before, at and after it, removing duplicate selections. Each monitor sees the same policy, tool definitions and conversation up to that decision. It sees no future messages, private labels or final database result. These sampled windows support local timing comparisons, not a claim about the first warning anywhere in a conversation. Across pilot and new-task study, 1,625 of the 1,626 planned monitor evaluations return valid scores; the remaining output stays missing.

To reduce interference from the simple interaction rule, we exclude entire conversations that violate it, along with policy--tool disagreements. The new-task study then supplies 47 matched decision points from seven tasks. Within each task we calculate the difference between the warning rate under the broad question and the rate under the six-rule question, then average the task differences equally. A positive difference means more warnings under the broad question.

\begin{table}[t]\centering\small
\caption{The same recorded decisions, two monitoring questions. Positive differences mean more warnings when asked about any policy violation; negative differences mean more warnings when asked about the six selected rules. These are changes in warning frequency, not proven changes in accuracy.}
\label{tab:substantive-targets}
\begin{tabular}{lrrr}\toprule
Monitor & Decision points & Tasks & Difference (percentage points)\\\midrule
Gemini 3.8 Flash & 47 & 7 & +45.2\\
GPT-5.4 Mini & 47 & 7 & $-9.5$\\
\bottomrule\end{tabular}
\end{table}

The two models move in opposite directions (Table~\ref{tab:substantive-targets}). Asking a narrower question therefore does not universally reduce warnings. Removing one task at a time leaves Gemini's difference between +36.1 and +52.8 percentage points, and GPT-5.4 Mini's between $-15.3$ and 0.0. The artifact also groups tasks by shared visible entities and reports sensitivity to those dependencies. The extension uses 2,000 task bootstrap resamples; these describe the small observed sample, not a deployment population.

A smaller comparison further requires a matching final database and a clean result on the applicable selected rules. It contains nine conversations from four tasks. Gemini warns on two under the broad question and none under the six-rule question; GPT-5.4 Mini warns on four and six, respectively. These counts are easy to misread: a broad warning might concern a policy we did not label. They measure warning frequency on a specified comparison group, not comprehensive false-alarm rates.

\paragraph{Observed example: changing the question changes the warning.}
In retail task 106 (GPT-5.4 Mini actor, repeat 0), message 32 proposes a T-shirt exchange on a delivered order after customer confirmation. The selected checks find the order status and same-product requirement satisfied; the tool accepts the exchange and the final database matches. This conversation has no interaction violation. On the identical prefix, Gemini scores the broad target at 1.0, citing a missing reminder to confirm all exchange items, but scores the six-rule target at 0.0. That reminder is outside the six checks, so this change cannot by itself be called improved accuracy. GPT-5.4 Mini moves in the opposite direction, from 0.12 to 0.86. Its latter explanation says no clear violation of the six constraints is evident despite the high score. We retain that inconsistent output: explanations do not override scores or supply ground-truth labels. This selected example explains how the comparison works; the full counts, rather than this illustration, describe the observed frequency.

The extension therefore supports a limited result: the requested target changes warnings even after removing conversations with interaction-rule violations. It does \emph{not} establish broad detection of substantive violations. The absence of positive examples in the strict subset, shared task entities, public training tasks and incomplete policy coverage remain material limitations. The next studies deliberately introduce pressure requests and permitted counterparts to observe intervention opportunities. That change supplies different evidence; it does not retroactively add positive examples to this strict subset.
